%% Region-Grounded Forced-Choice Attribute Verification for
%% Multimodal Hallucination Detection
%%
%% IEEE two-column, 6 pages. Build with:
%%     pdflatex main.tex && pdflatex main.tex
%%
%% Every number in this paper comes from results/tables/*.csv, which are
%% regenerated by `python reproduce.py`. Numbers marked TEST are from the
%% single test-split run (M5); all others are development-split.

\documentclass[conference]{IEEEtran}
\IEEEoverridecommandlockouts
\usepackage{cite}
\usepackage{amsmath,amssymb}
\usepackage{graphicx}
\usepackage{booktabs}
\usepackage{url}
\usepackage[hidelinks]{hyperref}

\begin{document}

\title{Region-Grounded Forced-Choice Attribute Verification\\
for Multimodal Hallucination Detection}

\author{\IEEEauthorblockN{Vanuj Gangrade, Aryan Bhandari, Harsh Parmeshwaram}
\IEEEauthorblockA{School of Computer Science and Engineering\\
Vellore Institute of Technology, Vellore, India\\
\{23BDS0282, 23BDS0223, 23BDS0007\}@vitstudent.ac.in}
\thanks{Guide: Dr.\ Gautam Majumdar, VIT Vellore.}}

\maketitle

\begin{abstract}
Vision--language models hallucinate attributes: they assert that an object has a
property the image does not support. Existing detectors check such claims by
scoring the claim text against the whole image with a contrastive model and
thresholding the score, an approach known to be unreliable because contrastive
models bind attributes to objects poorly. Two remedies are widely assumed to
help: cropping to the named object, and replacing the threshold with a forced
choice between competing attribute values. We implement both, and isolate them
with a $2\times2$ ablation on 3{,}858 development questions from AMBER,
confirmed on a held-out test split of 1{,}690 questions with all thresholds
frozen from the development fit. The combination improves accuracy by
$+0.2260$ (95\% CI $[+0.2001, +0.2514]$) on the test split, closely matching the
$+0.2232$ measured on development. The ablation, however, contradicts the common
assumption about \emph{why}: forced choice alone accounts for $+0.2036$ of that,
whereas region grounding alone accounts for $+0.0201$ --- an interval that clears
zero only marginally. A positive interaction observed on development
($+0.0223$) did \emph{not} replicate on test ($+0.0023$), and we report that
failure. We further show the gain does not depend on the benchmark supplying the
competing attribute: substituting a fixed antonym list written before any result
was computed, and scoring each question in isolation, still yields $+0.1901$ over
the baseline. Finally, we
document three construction artifacts in AMBER that allow trivial,
image-independent strategies to score $1.0000$, and show that the question types
supporting our result are free of them. The system uses two open-weight models
totalling 358M parameters, no language model, and no paid API, and runs in
1.65\,GiB of VRAM.
\end{abstract}

\begin{IEEEkeywords}
multimodal hallucination, vision--language models, attribute binding, ablation
study, benchmark artifacts
\end{IEEEkeywords}

\section{Introduction}

A vision--language model asked to describe a photograph may report that the sky
is \emph{gloomy} when it is plainly sunny. This is \emph{attribute
hallucination}: the object is present and correctly named, but a property
ascribed to it is unsupported. It is harder to detect than object
hallucination, because the object really is there and a detector cannot simply
check for its presence.

The standard detection recipe scores the claim text against the whole image with
a contrastive vision--language model and thresholds the resulting
score. Contrastive models are known to bind attributes to objects
poorly~\cite{yuksekgonul2023,thrush2022}, and UNIHD~\cite{chen2024} states the
consequence directly: it reports minimal improvement on attribute-level
hallucination and attributes this to the absence of a specialised attribute
tool.

Two changes are widely assumed to close that gap. \textbf{Region grounding}
localises the object named in the claim and evaluates the crop rather than the
whole image, on the intuition that the surrounding scene is what confuses the
binding. \textbf{Forced choice} replaces the absolute threshold with an argmax
over mutually exclusive candidate attributes, on the intuition that a relative
comparison is easier than an absolute judgement.

Both are plausible. Neither, to our knowledge, has been isolated against a
controlled baseline on the same data. This paper does that, and the result is not
the one the intuitions predict.

\textbf{Contributions.}
\begin{enumerate}
\item A $2\times2$ ablation isolating region grounding and forced choice
      independently, on 3{,}858 AMBER attribute questions, with confidence
      intervals bootstrapped over images.
\item The finding that \emph{forced choice carries almost the entire effect}.
      On the test split, region grounding alone contributes $+0.0201$ against
      forced choice's $+0.2036$. Two claims were withdrawn during this work on
      stronger evidence --- region grounding as null, and a superadditive
      interaction --- and we report both withdrawals rather than only the
      surviving numbers.
\item Two controls that rule out the obvious objections: a base-rate-matched
      threshold baseline, and a variant in which the competing attribute comes
      from a pre-registered external antonym list rather than the benchmark.
\item Three construction artifacts in AMBER that let image-independent
      strategies score $1.0000$, reported rather than exploited.
\end{enumerate}

\section{Related Work}

\textbf{Decompose-and-verify pipelines.} FaithScore, Woodpecker and
UNIHD~\cite{chen2024} decompose a generated caption into atomic claims and verify
each with an external tool. The pipeline shape is established and is not claimed
here as a contribution; we adopt it and specialise the attribute step.

\textbf{Region-level scoring.} ESREAL~\cite{esreal2024} performs region-level,
type-specific hallucination scoring. Region grounding as a technique is therefore
prior art. Our contribution is not the technique but its \emph{measurement}: we
quantify how much it contributes on its own, and find the answer to be small.

\textbf{Per-type reliability.} DHCP, HALP and CADMP report per-type breakdowns
establishing that attribute hallucination is harder than object hallucination.
That attributes are harder is not a finding of this work; it is our premise.

\textbf{Attribute binding.} Yuksekgonul et al.~\cite{yuksekgonul2023} and Thrush
et al.~\cite{thrush2022} show contrastive vision--language models behave close to
bag-of-words on compositional attribute tasks. Our forced-choice result is
consistent with their diagnosis: a relative comparison sidesteps the absolute
calibration these models lack.

\section{Method}

Given an image $I$ and a claim that object $o$ has attribute $a$, the system
emits \textsc{yes} or \textsc{no}.

\subsection{Region grounding}
An open-vocabulary detector (OWLv2-base~\cite{owlv2}) is prompted with
``a photo of a \{$o$\}''. Detections scoring at least $0.10$ are kept,
de-duplicated by non-maximum suppression at IoU $0.5$, and sorted by score. The
highest-scoring box is expanded by $10\%$ per side, clipped to the image, and
enforced to a minimum of $64\times64$ pixels. When no detection survives, the
crop falls back to the whole image and the record is flagged; we report the
fallback rate rather than hiding it.

\subsection{Forced choice}
A contrastive scorer (SigLIP-base~\cite{siglip}) scores the prompt
``a photo of a \{$a$\} \{$o$\}'' for each candidate attribute. Under
thresholding, each candidate is compared independently against a threshold
$\tau$. Under forced choice, the candidate with the higher score wins and the
other is answered \textsc{no}. Candidates are ordered alphabetically so option
position cannot encode the answer.

\subsection{The $2\times2$}
\begin{table}[t]
\caption{The ablation design. Cell A is the standard baseline; Cell D is the
proposed method.}
\label{tab:design}
\centering
\begin{tabular}{lcc}
\toprule
& \textbf{Threshold} & \textbf{Forced choice} \\
\midrule
\textbf{Whole image} & A (baseline) & B \\
\textbf{Cropped region} & C & D (proposed) \\
\bottomrule
\end{tabular}
\end{table}

Table~\ref{tab:design} gives the design. The prompt template is identical in all
four cells. $\tau$ is fitted on the development split for the threshold cells;
the forced-choice cells have \emph{no free parameter at all}.

\section{Experimental Setup}

\textbf{Data.} AMBER~\cite{amber2023}. Its \texttt{state} and \texttt{action}
questions occur in contrastive pairs over the same (image, object): one true
attribute and one false. We verified 2{,}378 clean state pairs and 396 action
pairs, so the forced-choice option set is supplied by the benchmark and no
antonym lexicon is needed on the evaluation path.

\textbf{Splits.} By image, never by question, so no image appears in both splits:
702 development images and 302 test images, frozen with a fixed seed. The test
split was opened exactly once, after all development work was complete.

\textbf{Metrics.} Accuracy over individual question ids, for comparability with
published AMBER numbers. Confidence intervals are bootstrapped over
\emph{images} (10{,}000 resamples), not questions, because questions from one
image are correlated. Comparisons between cells use a paired bootstrap on the
same resampled images.

\textbf{Models.} OWLv2-base-patch16-ensemble (155M) and SigLIP-base-patch16-224
(203M), both zero-shot, both pinned to exact revision hashes. Never resident
simultaneously: each loads, runs its stage, and is freed. Peak VRAM 1.65\,GiB.

\section{Results}

\subsection{The ablation}

\begin{table}[t]
\caption{Accuracy by cell. Development: 3{,}858 questions over 696 images, with
$\tau$ fitted on the evaluation data (which favours the baseline). Test:
1{,}690 questions over 296 images, with $\tau$ \emph{frozen} from the
development fit and nothing fitted on test.}
\label{tab:cells}
\centering
\begin{tabular}{lcc}
\toprule
\textbf{Cell} & \textbf{Dev} & \textbf{Test} \\
\midrule
A \;\; whole image + threshold & 0.5829 & 0.5751 \\
B \;\; whole image + forced choice & 0.7688 & 0.7787 \\
C \;\; cropped + threshold & 0.5980 & 0.5953 \\
\textbf{D \;\; cropped + forced choice} & \textbf{0.8061} & \textbf{0.8012} \\
\midrule
A$'$ base-rate matched & 0.5770 & 0.5775 \\
C$'$ base-rate matched & 0.5905 & 0.6012 \\
D-ext external contrast & 0.7621 & 0.7602 \\
Position-only (artifact) & 1.0000 & 1.0000 \\
\bottomrule
\end{tabular}
\end{table}

\begin{table}[t]
\caption{Attribution on the \textbf{test} split. Paired bootstrap over 296
images, 10{,}000 resamples. Development values in the last column for
comparison.}
\label{tab:attrib}
\centering
\small
\begin{tabular}{lccc}
\toprule
\textbf{Contrast (isolates)} & $\Delta$ & \textbf{95\% CI} & \textbf{Dev} \\
\midrule
B $-$ A \; forced choice & $+0.2036$ & $[+0.1781, +0.2288]$ & $+0.1858$ \\
B $-$ A$'$ \; net of base rate & $+0.2012$ & $[+0.1756, +0.2269]$ & $+0.1918$ \\
C $-$ A \; region grounding & $+0.0201$ & $[+0.0012, +0.0391]$ & $+0.0150$ \\
D $-$ C \; forced ch.\ $\mid$ crop & $+0.2059$ & $[+0.1808, +0.2306]$ & $+0.2081$ \\
D $-$ B \; crop $\mid$ forced ch. & $+0.0225$ & $[-0.0012, +0.0462]^{\ast}$ & $+0.0373$ \\
\midrule
\textbf{D $-$ A} \; \textbf{both} & $\mathbf{+0.2260}$ & $\mathbf{[+0.2001, +0.2514]}$ & $+0.2232$ \\
\midrule
Interaction & $+0.0023$ & --- & $+0.0223$ \\
\bottomrule
\multicolumn{4}{l}{\footnotesize $^{\ast}$ includes zero: this dev finding did not replicate.}
\end{tabular}
\end{table}

Tables~\ref{tab:cells} and~\ref{tab:attrib} give the result. On the held-out
test split, with every threshold frozen from the development fit, the combination
improves on the baseline by $+0.2260$, closely matching the $+0.2232$ measured on
development. \textbf{The headline effect is stable across splits}, and the
baseline received no threshold tuned on the data it was scored on.

The attribution is the substantive finding. \textbf{Forced choice alone accounts
for $+0.2036$ of the $+0.2260$} --- essentially all of it. Region grounding alone
accounts for $+0.0201$, an interval whose lower bound is $+0.0012$: it clears
zero, but only just.

\subsubsection*{Two withdrawn claims}
We report two claims that this work made and then retracted on stronger evidence.

First, on an early 100-pair sample we reported region grounding as \emph{null}
($+0.0100$, CI $[-0.0490, +0.0680]$). The full development set narrowed the
interval until it cleared zero, and we withdrew that claim.

Second, on the development set we found a positive interaction: cropping was
worth a further $+0.0373$ once forced choice was in place
(CI $[+0.0217, +0.0529]$), with superadditivity of $+0.0223$. \textbf{This did
not replicate.} On test, D $-$ B is $+0.0225$ with CI $[-0.0012, +0.0462]$, which
includes zero, and the interaction falls to $+0.0023$. On the test split the two
changes behave additively, with forced choice supplying almost the entire effect.

We report both withdrawals because a claim stated, tested and retracted is
evidence about the process that produced the numbers we do keep. The surviving
description is narrow: \emph{the decision rule is what matters}; region grounding
contributes a small, marginal and split-dependent amount, and should not be
presented either as null or as the mechanism.

\subsection{Control 1: the base rate}

Forced choice emits exactly one \textsc{yes} per contrastive pair. Since AMBER's
attribute pairs are exactly balanced, this hands it a correct answer distribution
for free, which the threshold cells do not receive: Cell A answers \textsc{yes}
to $51.4\%$ of questions and Cell C to $37.5\%$, against a true rate of $50\%$.

Cells A$'$ and C$'$ remove that asymmetry by choosing $\tau$ to match the known
base rate instead of to maximise accuracy. \emph{Giving the baseline the prior
widens the gap rather than closing it}: forced choice is worth
$(+0.2036, +0.2012)$ on the whole image and $(+0.2059, +0.2000)$ on the crop
(test split), quoting the raw and base-rate-matched contrasts as a pair. Matching the base rate
costs the threshold cells accuracy, because their accuracy-maximising threshold
was deliberately unbalanced.

Note the two constraints differ. A$'$/C$'$ impose a \emph{global} base rate;
forced choice imposes a \emph{pairwise} one. The latter is strictly stronger and
is unavailable to a thresholding rule, which has no representation of the pair.
Taking it requires comparing the candidates against each other --- which is
forced choice. The remaining margin is therefore the contribution, not an
unremoved confound.

\subsection{Control 2: an external contrast}

A reviewer may object that the method exploits pair structure the baseline cannot
access, since AMBER supplies the competing attribute. We remove that dependency
entirely.

Cell D-ext draws its competitor from a fixed antonym list, written from the
ordinary meaning of each word and \emph{committed to version control before any
D-ext result was computed}. Each question is answered in isolation: ``Is the sky
sunny?'' is scored against the list's opposite of \emph{sunny}, and its paired
question against the opposite of \emph{gloomy}. The two never interact, so D-ext
may answer \textsc{yes} to both halves of a pair --- and does, on 17\% of pairs.
It therefore surrenders the pairing advantage completely.

The list covers $94.0\%$ of test questions and disagrees with AMBER's own
negative on $44.1\%$ of pairs. On the matched covered subset of the test split,
D-ext scores $0.7602$ against the baseline's $0.5702$: \textbf{$+0.1901$, CI
$[+0.1645, +0.2157]$}. It trails full Cell D by $-0.0428$, the measured price of
giving up the pairing advantage. Restricted further to questions where the external word \emph{differed}
from AMBER's, D-ext still beats the baseline. The gain does not depend on the
benchmark's own contrast.

\section{Benchmark Artifacts in AMBER}

While building the supporting modules we found three independent construction
artifacts, each allowing a strategy that never opens an image to score $1.0000$.

\begin{enumerate}
\item \textbf{Attribute pair id ordering.} In all 2{,}774 clean attribute pairs,
      the true attribute holds the \emph{lower} question id, and the two ids are
      always adjacent. Answering \textsc{yes} to the lower id of each pair scores
      $1.0000$.
\item \textbf{All-negative existence gold.} All 4{,}924
      \texttt{discriminative-hallucination} questions ask about absent objects;
      every gold answer is \textsc{no}. Answering \textsc{no} always scores
      $1.0000$, and accuracy is not an interpretable metric for this subset.
\item \textbf{Per-type constant relation gold.} \texttt{discriminative-relation}
      is entirely \textsc{yes} (975 questions) and \texttt{relation} entirely
      \textsc{no} (689).
\end{enumerate}

We report these rather than exploit them. For artifact~1 we implement the
exploiting baseline explicitly, label it, and verify by permutation that no cell
in our system reads question ids: under permuted ids the exploiting baseline
collapses to $0.41$ while Cells A and D are bit-identical.

Critically, \textbf{the three question types our result rests on ---
\texttt{state}, \texttt{action} and \texttt{number} --- are exactly balanced and
free of these artifacts.} All three degenerate types fall in the supporting-module
scope, for which we claim no novelty.

\section{Full Pipeline}

The attribute module is deployed alongside existence, counting and relation
checkers so the system is complete. Consistent with the artifacts above,
existence is reported by \emph{false-positive rate} rather than accuracy, and
relation is reported per type and pooled, never pooled alone. Counting is
unaffected: its labels are balanced, so accuracy is meaningful. Full figures are
in Table~\ref{tab:pipeline}.

Two of these deserve comment. Counting answers yes/no correctly $72\%$ of the
time while producing the \emph{correct count} only $36\%$ of the time (mean
absolute error $1.15$ objects): because roughly half the questions are gold-\textsc{no},
a miscount often yields the right answer for the wrong reason. And the relation
module barely exceeds its trivial baseline ($0.6056$ against $0.5851$). Both are
off-the-shelf components included for system completeness, not contributions, and
we report them as they stand.

\begin{table}[t]
\caption{Full pipeline, development split. Metrics differ by module because three
question types have degenerate labels.}
\label{tab:pipeline}
\centering
\small
\begin{tabular}{llcc}
\toprule
\textbf{Module} & \textbf{Metric} & \textbf{Value} & \textbf{Trivial} \\
\midrule
Attribute (Cell D) & accuracy & 0.8061 & 0.5000 \\
Counting & accuracy & 0.7196 & 0.5000 \\
Counting & exact count & 0.3632 & --- \\
Existence & false-pos.\ rate & 0.1244 & 1.0000$^{\dagger}$ \\
Relation (pooled) & accuracy & 0.6056 & 0.5851 \\
\bottomrule
\multicolumn{4}{l}{\footnotesize $^{\dagger}$ degenerate labels; accuracy not interpretable.}
\end{tabular}
\end{table}

\section{Cost}

The system uses 358M parameters across two open-weight models, no language model,
and no paid API. Inference cost per question is dominated by detection: $669$\,ms per counting
question and $160$\,ms per relation question end to end, against $32$\,ms for an
attribute question once detections are cached. Peak VRAM during two-model
staging was $1.65$\,GiB of the $4.00$\,GiB available on the target laptop GPU,
achieved by never holding both models resident. Monetary API cost is exactly
zero.

\section{Limitations}

\textbf{Development-split thresholds are fitted on the evaluation data.} Cells A,
C, A$'$ and C$'$ each fit $\tau$ on the questions they are scored on. This
favours the \emph{baseline}, so the reported comparison is conservative; the
headline test-split result uses $\tau$ frozen from the development fit.

\textbf{Effective sample is smaller than $n$.} The 3{,}858 questions rest on 661
distinct (object, positive, negative) triples, and the bootstrap resamples 696
images.

\textbf{Detector fallback.} On $3.4\%$ of questions no object was detected and
the crop was the whole image, making Cell D behave as Cell B there.

\textbf{Region grounding's effect is marginal and may not be robust.} Its test
interval clears zero by $+0.0012$, and the development interaction it participated
in did not replicate. A further split could plausibly place it either side of
zero.

\textbf{Forced choice requires a competing hypothesis.} Where the benchmark
supplies one, Cell D applies. Where it does not, an external contrast is needed,
and for roughly a fifth of action attributes (\emph{swim}, \emph{jump},
\emph{surf}) no natural opposite exists in English. This is a scope condition of
the method, not a defect of the experiment.

\textbf{Confidences are uncalibrated} and calibration is not a claim of this work.

\textbf{Scope.} English only; AMBER's closed 340-object vocabulary; zero-shot
with no training or fine-tuning; no OCR or scene text.

\section{Conclusion}

Region-grounded forced-choice attribute verification improves attribute
hallucination detection by $+0.2260$ over the standard thresholding baseline on a
held-out test split, with thresholds frozen from development. The ablation shows
this is overwhelmingly attributable to the \emph{decision rule}, not to region
grounding: forced choice alone is worth $+0.2036$ and cropping alone $+0.0201$,
and the positive interaction we measured on development did not replicate. The
gain survives removing the benchmark's own contrast. Along the way we document three construction
artifacts in AMBER that admit perfect image-independent strategies, and show that
the question types supporting our result are free of them.

The practical implication is narrow and useful: when adding an attribute tool to
a decompose-and-verify pipeline, the change that pays is asking the model to
choose between competing attributes rather than to judge one in isolation.
Localisation helps, but far less than is commonly assumed --- and by an amount
small enough that we would not rely on it without further replication.

\begin{thebibliography}{9}
\bibitem{yuksekgonul2023} M.~Yuksekgonul et al., ``When and why vision-language
models behave like bags-of-words, and what to do about it?'' \emph{ICLR}, 2023.
\bibitem{thrush2022} T.~Thrush et al., ``Winoground: Probing vision and language
models for visio-linguistic compositionality,'' \emph{CVPR}, 2022.
\bibitem{chen2024} X.~Chen et al., ``Unified hallucination detection for
multimodal large language models,'' 2024.
\bibitem{amber2023} J.~Wang et al., ``AMBER: An LLM-free multi-dimensional
benchmark for MLLMs hallucination evaluation,'' 2023.
\bibitem{esreal2024} M.~Kim et al., ``ESREAL: Exploiting semantic reconstruction
to mitigate hallucinations in vision-language models,'' \emph{ECCV}, 2024.
\bibitem{owlv2} M.~Minderer, A.~Gritsenko, and N.~Houlsby, ``Scaling
open-vocabulary object detection,'' \emph{NeurIPS}, 2023.
\bibitem{siglip} X.~Zhai et al., ``Sigmoid loss for language image
pre-training,'' \emph{ICCV}, 2023.
\end{thebibliography}

\end{document}
